\scanpage{76}
\mspage{82}

et il suffit de prouver que cette flèche est un iso. Pour ceci, on regarde fibre par fibre, et on utilise le fait que ``\deleted{\unsure{les}} les images directes se calculent fibre par fibre''.

\uline{Question} Soit $f:X\longrightarrow Y$ une $\mathbb D$-fibration, et soit $F\in\mathrm{Ob}\,\mathbb D(Y)$. Alors $f_* f^*F$ se calcule-t-il \deleted{\unsure{fibre}} \unsure{p. r.|par} \deleted{\unsure{par|pro}} \inserted{aux} fibres ? Même question pour $f_! f^*F$. Pour la première propriété, concernant $f_*$, on montre que cela marche si $F$ est de la forme $p_*(\deleted{\unsure{N}}M_{Y'})$, où $p:Y'\longrightarrow Y$ est $\mathbb D$-\uline{propre}, et $M\in\mathrm{Ob}\,A$.

D'autre part, la sous-catégorie pleine \inserted{$C$} de $\mathbb D(Y)$ formée des $F$ pour lesquels ça marche est \inserted{(sauf erreur)} ``stable par $\mathbb D$-$\unsure{\mathrm{hot}|\mathrm{ho}}$-$\underset{\longleftarrow}{\mathrm{lim}}$'', en le sens suivant : si on a une catégorie $I$ avec objet final $e$ formant une sous-catégorie fermée (i.e. $e\longrightarrow x\ \Longrightarrow\ x=e$), et si on a $\mathscr F\in\deleted{\unsure{A}}\mathbb D(Y\times I)$, d'où
\[
\begin{aligned}
I^\circ&\longrightarrow\mathbb D(Y)\\
i&\longrightarrow\mathscr F_i
\end{aligned}
\]
(diagramme squelettique), si les $\mathscr F_i$ pour $i\ne e$ sont dans $C$, et si $\mathscr F$ est dans l'image essentielle de $j_*$, où $j:Y\times(I\setminus\{e\})\unsure{\longleftarrow|\longrightarrow}Y\times I$ est l'inclusion, alors $\mathscr F_e$ est aussi dans $C$.

On peut cependant considérer que si on veut que $f:X\longrightarrow Y$ satisfasse à toutes les propriétés coh. et hom. qui s'attachent à la notion d'une ``fibration'' (``loc\unsure{.} triviale''), on devrait mettre dans les conditions, que $f$ est \uline{lisse}, et \uline{propre}, donc que le chgt. de base par $f$ commute aux foncteurs $g_*,g_!$.
% archive sheet number: 75
